\chapter{Metodología de desarrollo}\label{chap:metodologia}
\label{sec:metodologia}

Esta sección describe el proceso con el que se desarrolló el sistema: el enfoque adoptado y su justificación, la planificación inicial y sus desvíos, los incrementos efectivamente ejecutados con sus fechas reales, la herramienta de seguimiento, los ajustes que surgieron de los resultados del modelo y el \emph{pipeline} de entrenamiento.

\section{Enfoque y justificación}
El proyecto se desarrolla mediante un proceso iterativo e incremental \parencite{Larman2003},
organizado en torno a casos de uso. Cada incremento toma un conjunto de casos de uso, lo lleva a través de
análisis, diseño, implementación y prueba, y entrega una versión ejecutable del sistema completo
sobre la que se planifica el incremento siguiente. Se elige este enfoque por tres motivos. Primero,
el comportamiento del modelo de detección no puede conocerse antes de entrenarlo y probarlo
contra la telemetría real del agente, de modo que los requerimientos se ajustan a partir de los
resultados de cada iteración. Segundo, el sistema tiene tres componentes desplegados en forma
independiente (agente, API y panel), y un incremento vertical que atraviesa los tres permite
verificar la integración temprano. Tercero, un equipo de dos personas no justifica los roles y
ceremonias de un marco como Scrum, pero sí la entrega frecuente de versiones verificables.

\section{Planificación inicial y desvíos}
\label{sec:plan-contra-real}
La planificación inicial, elaborada el 13 de junio de 2026 como diagrama de Gantt, organizaba el trabajo en fases por componente (infraestructura y \emph{pipeline} de datos, agente e integración con la API, iteración del modelo, API central, panel web, validación del puntaje y documentación), con hitos en las entregas del 25\,\%, 50\,\% y 75\,\% y la defensa final. La tabla~\ref{tab:plan-contra-real} compara cada fase con su ejecución real. Las fechas reales se toman del historial de versiones: el inicio es el primer \emph{commit} que nombra el trabajo y el fin, el último, por lo que son fechas de \emph{commit} y no de calendario de trabajo.

\begin{longtable}{>{\raggedright\arraybackslash}p{3.6cm}>{\raggedright\arraybackslash}p{2.0cm}>{\raggedright\arraybackslash}p{2.0cm}>{\raggedright\arraybackslash}p{6.4cm}}
	\caption{Planificación inicial contra ejecución real}
	\label{tab:plan-contra-real}\\
	\hline
	\textbf{Fase planificada} & \textbf{Plan} & \textbf{Real} & \textbf{Desvío y causa} \\
	\hline
	\endfirsthead
	\hline
	\textbf{Fase planificada} & \textbf{Plan} & \textbf{Real} & \textbf{Desvío y causa} \\
	\hline
	\endhead
	Infraestructura base y \emph{pipeline} de datos y modelo & 14/06--18/08 & 03/05--04/08 & Se adelantó: la preparación de CLUE-LDS y los primeros modelos empezaron antes de formalizar el plan. \\
	\hline
	Entorno de pruebas y datos sintéticos & 21/06--18/07 & -- & Se reemplazó por escenarios construidos sobre historiales reales, que luego también se retiraron (sección~\ref{sec:revision-metodo}). \\
	\hline
	Agente local e integración con la API & 19/07--18/08 & 17/07--12/08 & En plazo; cerró con la demostración del 50\,\%. \\
	\hline
	Panel web e integración & 19/08--20/10 & 17/07--27/08 & Se adelantó para que el incremento 1 cubriera CU-02 completo. \\
	\hline
	Modelo: iteración y reentrenamiento & 08/08--27/09 & 20/08--25/08 y 04/10--13/11 & El reentrenamiento (CU-03) se cerró en plazo; la revisión del modelo se extendió siete semanas más, como un incremento propio. \\
	\hline
	API central: desarrollo avanzado & 08/08--27/09 & 31/08--20/10 & Se extendió por la integración IAM y el endurecimiento, que no estaban en el plan. \\
	\hline
	Validación y calibración del puntaje & 19/08--05/10 & 04/10--13/11 & Se postergó y se rehízo con datos reales, criterios fijados antes de medir y una partición DEV/HOLDOUT fija. \\
	\hline
	Despliegue e infraestructura & 06/09--12/10 & 10/09--14/11 & Se extendió por las sucesivas versiones desplegadas en la nube. \\
	\hline
	Hito del 75\,\% & 20/10 & 20/10 & Sin desvío. \\
	\hline
	Cierre y defensa & 21/10--14/12 & desde 13/11 & En curso; incluye las correcciones de la entrega del 90\,\%. \\
	\hline
\end{longtable}
{\small Fuente: elaboración propia a partir de la planificación inicial y del historial del repositorio.\par}

El historial registra el primer \emph{commit} el 16 de abril de 2026 y la siguiente cantidad de \emph{commits} por mes: 4 en abril, 3 en mayo, 17 en julio, 119 en agosto, 26 en septiembre, 359 en octubre y 46 en noviembre (hasta el 16). La concentración en agosto y octubre coincide con los dos incrementos de producto y con la revisión de modelos.

\section{Incrementos ejecutados}
\label{sec:incrementos}
El plan por fases se ejecutó, en la práctica, como una secuencia de incrementos verticales. La tabla~\ref{tab:incrementos} los presenta con sus fechas reales, su alcance y la versión o hito con que cerró cada uno.

\begin{longtable}{>{\raggedright\arraybackslash}p{3.2cm}>{\raggedright\arraybackslash}p{2.4cm}>{\raggedright\arraybackslash}p{5.0cm}>{\raggedright\arraybackslash}p{3.4cm}}
	\caption{Incrementos ejecutados}
	\label{tab:incrementos}\\
	\hline
	\textbf{Incremento} & \textbf{Fechas reales} & \textbf{Alcance} & \textbf{Versión o hito} \\
	\hline
	\endfirsthead
	\hline
	\textbf{Incremento} & \textbf{Fechas reales} & \textbf{Alcance} & \textbf{Versión o hito} \\
	\hline
	\endhead
	0. Datos y modelos base & 03/05--04/08 & Fases 1 a 6 sobre CLUE-LDS: limpieza, agregado usuario-día, secuencias, Isolation Forest y LSTM Autoencoder & Artefactos de la versión~1 (02/09) \\
	\hline
	1. Detección y ajuste de políticas (CU-01 y CU-02) & 17/07--27/08 & API con puntaje híbrido, agente Windows con mitigación local, panel con umbrales y acciones editables y niveles configurables (RF-13) & Demostración del 50\,\% (12/08) \\
	\hline
	2. Reentrenamiento, alta de usuarios, IAM y endurecimiento (CU-03 y CU-04) & 20/08--20/10 & Captura de eventos crudos, planificador, criterio de promoción y recarga sin reinicio (RF-19 a RF-22); período de aprendizaje configurable (RF-24 y RF-25); adaptador IAM con Keycloak y Entra ID; resiliencia, deduplicación, clave por agente, cadena de \emph{hashes} y sesión del panel & Primera corrida automática completa (25/08); entrega del 75\,\% (20/10) \\
	\hline
	Despliegue en la nube (transversal) & 10/09--14/11 & Terraform, instancia, \emph{proxy} inverso y versiones publicadas & Versiones del 08/11, 14/11 y 16/11 \\
	\hline
	3. Revisión de modelos & 04/10--13/11 & Partición DEV/HOLDOUT fija, criterio de no inferioridad, retiro de escenarios construidos, tercer detector y etiquetado de triage (sección~\ref{sec:revision-modelos}) & Etiqueta \texttt{pre-limpieza-entrega} (13/11) \\
	\hline
	4. Cierre y correcciones del 90\,\% & 13/11--16/11 & Regresión del puntaje, pruebas del panel, cambios de modelo de la sección~\ref{sec:que-cambiar-primero} y evaluación única en HOLDOUT & Versión del 16/11 \\
	\hline
\end{longtable}
{\small Fuente: elaboración propia a partir del historial del repositorio.\par}

El primer incremento cubre la detección y la mitigación en el equipo (CU-01) y el ajuste de políticas desde el panel (CU-02), y se verificó con la demostración del 50\,\%. El segundo incorpora el reentrenamiento periódico del modelo (CU-03), el alta de usuarios nuevos con período de aprendizaje (CU-04), la integración con sistemas de gestión de identidad y el endurecimiento de la ingesta, la autenticación y el canal; se verificó con la primera corrida automática del reentrenamiento y con la entrega del 75\,\%. La revisión de modelos funcionó como un tercer incremento: tuvo un objetivo propio, criterios de decisión fijados antes de medir y una versión desplegada como resultado.

\section{Herramienta de seguimiento}
El seguimiento se hizo exclusivamente con el control de versiones: no se usaron tableros, gestores de incidencias ni planillas. Las ramas por funcionalidad (con prefijos como \texttt{feature/}, \texttt{analysis/} y \texttt{fix/}) hicieron las veces de tablero, y cada sesión de trabajo dejó en el repositorio un documento que registra qué se decidió y por qué. La revisión de cada cambio antes de integrarlo y un \emph{pipeline} de integración continua que ejecuta el análisis estático, la compilación y las pruebas automatizadas (sección~\ref{sec:integracion-continua}) completan el control. La división de tareas responde a la especialización de cada integrante: uno concentra el entrenamiento del modelo y el otro el agente, la API y el panel.

Como limitación, la historia es casi lineal (diez \emph{merges}, en su mayoría de la integración de la entrega del 75\,\%), por lo que las ramas no permiten reconstruir por sí solas la duración de cada tarea, y las fechas de la tabla~\ref{tab:incrementos} son las del primer y el último \emph{commit} de cada trabajo.

\section{Ajustes derivados de los resultados del modelo}
\label{sec:ajustes-modelo}
Varias decisiones de diseño surgieron de resultados medidos y no del plan inicial. La tabla~\ref{tab:ajustes-modelo} los resume con su fecha.

\begin{longtable}{>{\raggedright\arraybackslash}p{4.4cm}>{\raggedright\arraybackslash}p{7.6cm}>{\raggedright\arraybackslash}p{1.6cm}}
	\caption{Ajustes derivados de los resultados del modelo}
	\label{tab:ajustes-modelo}\\
	\hline
	\textbf{Ajuste} & \textbf{Resultado que lo motivó} & \textbf{Fecha} \\
	\hline
	\endfirsthead
	\hline
	\textbf{Ajuste} & \textbf{Resultado que lo motivó} & \textbf{Fecha} \\
	\hline
	\endhead
	Umbral de aprendizaje configurable & Distribución del tiempo que tardan los usuarios reales en acumular 20 eventos & 27/08 \\
	\hline
	Criterio de promoción por no inferioridad & El criterio anterior (``iguala o mejora'') decidió dos rechazos por ruido & 10/10 \\
	\hline
	Corte temporal fijo y \emph{pipeline} determinista & Con un corte dinámico, el criterio comparaba particiones distintas & 28/10 \\
	\hline
	Clasificación única de usuarios & Dos fases del \emph{pipeline} clasificaban distinto a los mismos usuarios & 28/10 \\
	\hline
	Evaluación solo con datos reales; retiro de los escenarios S1--S7 & Los escenarios construidos no representan ataques reales y orientaban las decisiones & 29/10 \\
	\hline
	LSTM fuera del camino por defecto & Permutar la secuencia no cambia su ordenamiento (Spearman 0,989) y el 95,4\,\% de la varianza de su riesgo está entre usuarios & 30/10 \\
	\hline
	Detección de novedad como tercer detector & Criterio fijado sobre CLUE-LDS DEV & 30/10 \\
	\hline
	Markov fuera del camino por defecto & Su regla de piso explicaba 177 de 189 días con oscilación del nivel & 31/10 \\
	\hline
	LANL 2015 retirado como evidencia & Una sola intrusión, casi toda desde un equipo \parencite{Larroche2026} & 31/10 \\
	\hline
	Regresión contra la simulación del camino por defecto & El oráculo anterior estaba desactualizado y contaba eventos duplicados & 16/11 \\
	\hline
	Novedad solo informativa; país fuera del modelo con una regla explícita de ubicación & Comparación de opciones con un criterio fijado antes de medir (sección~\ref{sec:que-cambiar-primero}) & 16/11 \\
	\hline
\end{longtable}
{\small Fuente: elaboración propia a partir del historial del repositorio.\par}

\section{\emph{Pipeline} de entrenamiento}
El \emph{pipeline} de entrenamiento del modelo va desde la preparación del conjunto de datos
hasta la exportación de los artefactos que consume la API, y se ejecuta en forma separada del
ciclo de desarrollo del software. En el camino por defecto, sus fases son el análisis
exploratorio y el perfilado del conjunto de datos; la limpieza y la clasificación de los usuarios
por volumen y tipo; la ingeniería de variables; el entrenamiento de Isolation Forest; y la
configuración del puntaje, que fija la vida media del decaimiento y los umbrales y exporta las
referencias y la calibración, junto con el perfil de Markov que usa el camino alternativo
\texttt{markov}. El camino alternativo \texttt{lstm} agrega el
entrenamiento del LSTM Autoencoder, una auditoría bloqueante de consistencia de los datos y la
combinación de ambos modelos (sección~\ref{sec:caminos-alternativos}). El perfilado de completitud de los campos, que habría anticipado la falta casi total del país, se hizo después de elegir las variables (sección~\ref{sec:ea-datasets}).
